\documentclass[10pt]{article}
\providecommand{\FIGDIR}{fig}
\newcommand{\DIGESTDATE}{31 August 2026}
% ===================== PM Research Watch - shared preamble =====================
\usepackage[T1]{fontenc}
\usepackage[utf8]{inputenc}
\usepackage{amsmath,amssymb}
\usepackage{tgpagella}
\usepackage{tgheros}
\renewcommand{\sfdefault}{qhv}
\renewcommand{\ttdefault}{lmtt}
\usepackage{microtype}
\usepackage[english]{babel}

\usepackage[a4paper,top=17mm,bottom=17mm,left=17mm,right=17mm,
            headsep=6mm,footskip=10mm]{geometry}
\usepackage{graphicx}
\usepackage{xcolor}
\usepackage{colortbl}
\usepackage{tikz}
\usetikzlibrary{calc}
\usepackage[most]{tcolorbox}
\usepackage{titlesec}
\usepackage{fancyhdr}
\usepackage{booktabs}
\usepackage{array}
\usepackage{longtable}
\usepackage{enumitem}
\usepackage{multicol}
\usepackage{ragged2e}
\usepackage[hidelinks]{hyperref}

% ------------------------------------------------------------------ palette
\definecolor{Ink}    {HTML}{14262E}
\definecolor{Deep}   {HTML}{0F4C5C}
\definecolor{Teal}   {HTML}{1F7A8C}
\definecolor{Sky}    {HTML}{4FA3C4}
\definecolor{Sage}   {HTML}{6FA287}
\definecolor{Amber}  {HTML}{E3A008}
\definecolor{Clay}   {HTML}{B5643C}
\definecolor{Coral}  {HTML}{D1495B}
\definecolor{Violet} {HTML}{6A5B9E}
\definecolor{Slate}  {HTML}{4A6572}
\definecolor{Paper}  {HTML}{FBF9F5}
\definecolor{Mist}   {HTML}{EEEAE2}
\definecolor{Rule}   {HTML}{D9D3C9}

\pagecolor{Paper}
\color{Ink}

\hypersetup{colorlinks=true, linkcolor=Deep, urlcolor=Teal, citecolor=Teal,
            pdftitle={PM Research Watch}, pdfauthor={Automated daily digest}}

% ------------------------------------------------------------------ headings
\titleformat{\section}
  {\sffamily\bfseries\large\color{Deep}}{}{0pt}
  {}[{\vspace{1.2mm}\color{Amber}\titlerule[1.1pt]}]
\titlespacing*{\section}{0pt}{5.5mm}{2.8mm}

\titleformat{\subsection}
  {\sffamily\bfseries\normalsize\color{Teal}}{}{0pt}{}
\titlespacing*{\subsection}{0pt}{3.6mm}{1.4mm}

\setlist[itemize]{leftmargin=4.2mm, itemsep=1.1mm, topsep=1.3mm, parsep=0pt,
                  label=\textcolor{Teal}{\footnotesize$\blacktriangleright$}}

% ------------------------------------------------------------------ header/footer
\pagestyle{fancy}
\fancyhf{}
\renewcommand{\headrulewidth}{0pt}
\renewcommand{\footrulewidth}{0pt}
\fancyhead[L]{\sffamily\scriptsize\color{Slate}PM RESEARCH WATCH\;\textcolor{Rule}{\textbar}\;\DIGESTDATE}
\fancyhead[R]{\sffamily\scriptsize\color{Slate}Particulate matter monitoring \& health effects}
\fancyfoot[C]{\sffamily\scriptsize\color{Slate}\thepage}

% ------------------------------------------------------------------ components
% metric tile
\newcommand{\metric}[3]{%
  \begin{tcolorbox}[enhanced, colback=white, colframe=Rule, boxrule=0.5pt,
    arc=1.6mm, left=2.2mm, right=2.2mm, top=1.8mm, bottom=1.8mm,
    borderline west={1.7mm}{0pt}{#3}, nobeforeafter, valign=center,
    height=17mm, width=\linewidth]
    \RaggedRight
    {\sffamily\bfseries\fontsize{17}{17}\selectfont\color{#3}#1}\\[0.5mm]
    {\sffamily\fontsize{6.5}{7.4}\selectfont\color{Slate}#2}
  \end{tcolorbox}}

% pull-quote / key insight box
\newtcolorbox{keybox}[1][]{enhanced, colback=Mist, colframe=Mist,
  arc=1.6mm, left=3.2mm, right=3.2mm, top=2.4mm, bottom=2.4mm,
  borderline west={1.7mm}{0pt}{Amber}, #1}

% subtopic band
\newcommand{\band}[2]{%
  \vspace{3.2mm}\noindent
  \begin{tcolorbox}[enhanced, colback=#2, colframe=#2, arc=1.2mm,
    left=2.6mm, right=2.6mm, top=1.1mm, bottom=1.1mm, boxrule=0pt,
    width=\linewidth, nobeforeafter]
    {\sffamily\bfseries\small\color{white}#1}
  \end{tcolorbox}\vspace{1.6mm}}

% one paper entry: {short cite}{title}{journal}{doi}{body}
\newcommand{\paperentry}[5]{%
  \par\vspace{1.9mm}\noindent
  {\sffamily\bfseries\footnotesize\color{Ink}#1}\;%
  {\sffamily\fontsize{7}{8}\selectfont\color{Slate}#3}\\[0.4mm]
  {\footnotesize\itshape\color{Deep}#2}\\[0.6mm]
  {\footnotesize #4}\\[0.3mm]
  {\fontsize{6.6}{7.6}\selectfont\color{Slate}\texttt{doi:}\,\href{https://doi.org/#5}{\texttt{\detokenize{#5}}}}
  \par}

% caption style
\newcommand{\figcap}[1]{\par\vspace{1.0mm}
  {\fontsize{7.2}{8.6}\selectfont\color{Slate}\textsf{\textbf{Fig.}}\; #1\par}\vspace{1.5mm}}

\providecommand{\permil}{\textperthousand}

\setlength{\parindent}{0pt}
\setlength{\parskip}{0pt}
\renewcommand{\arraystretch}{1.18}


\begin{document}

% ==================================================================== MASTHEAD
\thispagestyle{empty}
\begin{tikzpicture}[remember picture, overlay]
  \fill[Deep] (current page.north west) rectangle
      ($(current page.north east)+(0,-67mm)$);
  \fill[Amber] ($(current page.north west)+(0,-67mm)$) rectangle
      ($(current page.north east)+(0,-68.4mm)$);
  \foreach \x/\y/\r in {22/12/0.9, 41/26/1.5, 63/9/0.7, 88/31/1.1, 112/17/1.9,
                        138/29/0.8, 159/13/1.3, 178/54/1.0, 196/21/1.6,
                        30/44/0.6, 74/57/1.2, 125/48/0.9, 168/6/0.8, 191/59/0.7,
                        52/18/1.0, 100/60/0.6, 148/20/1.1, 66/49/1.4, 205/45/0.8}
    \fill[Sky, opacity=0.30] ($(current page.north west)+(\x mm,-\y mm)$) circle (\r mm);
\end{tikzpicture}

\vspace*{4mm}
{\sffamily\fontsize{7.6}{9}\selectfont\color{Sky}%
 AUTOMATED DAILY LITERATURE DIGEST\;\textbullet\;PhD RESEARCH WATCH\par}
\vspace{2.2mm}
{\sffamily\bfseries\fontsize{27}{29}\selectfont\color{white}%
 Particulate Matter\par}
{\sffamily\fontsize{27}{29}\selectfont\color{Sky}%
 Monitoring \& Health Effects\par}
\vspace{5.0mm}
{\sffamily\fontsize{8.4}{10}\selectfont\color{white}%
 Issue 2026-08-31\;\;\textcolor{Amber}{\textbar}\;\;Entry window 31 August 2026 (single day)\;\;%
 \textcolor{Amber}{\textbar}\;\;16 records in scope, 31 screened out\par}
\vspace{18mm}

% -------------------------------------------------------------- metric strip
\noindent
\begin{minipage}[t]{0.190\linewidth}\metric{16}{RECORDS IN SCOPE\\(31 SCREENED OUT)}{Deep}\end{minipage}\hfill
\begin{minipage}[t]{0.190\linewidth}\metric{7}{SUBTOPIC CLUSTERS\\OF TEN OCCUPIED}{Teal}\end{minipage}\hfill
\begin{minipage}[t]{0.190\linewidth}\metric{5}{EXPOSURE ASSESSMENT\\\& MODELLING}{Amber}\end{minipage}\hfill
\begin{minipage}[t]{0.190\linewidth}\metric{2}{EFFECT ESTIMATES\\(ONE SOURCE)}{Clay}\end{minipage}\hfill
\begin{minipage}[t]{0.190\linewidth}\metric{10}{RECORDS WITH NO\\HEALTH ENDPOINT}{Sky}\end{minipage}

\vspace{6mm}

% -------------------------------------------------------------- signal of day
\section*{\sffamily\bfseries\color{Deep}Signal of the day}
\vspace{-1.5mm}{\color{Amber}\rule{\linewidth}{1.1pt}}\vspace{2.5mm}

\begin{keybox}
{\sffamily\bfseries\small\color{Deep}Two \textit{ES\&T} papers published on the same day
argue, from opposite directions, that PM$_{2.5}$ \emph{mass} is the wrong control
variable --- one on toxicity weighting, one on the exposure surfaces that feed it.}\\[1.4mm]
\footnotesize
Song et al.\ make the policy case: age-standardised PM$_{2.5}$-attributable death rates
have diverged worldwide over the past decade even as mass concentrations fell, which is
prima facie evidence against the uniform-toxicity assumption written into every current
standard. Weighting exposure by source-specific toxicity --- incomplete combustion,
especially residential solid fuel and wildfire, carrying the highest toxicity per unit
mass --- puts \textbf{24.0\% of the global toxicity-adjusted exposure burden on the 8.2\%
of the population that is low-income}. Li Q et al.\ supply the other half of the machinery
for China: a 1\,km daily PM$_{2.5}$ surface for 2013--2023 (cross-validated
$R^2$ 0.88 spatially, 0.82 temporally) fused with meta-analysed short-term effect
estimates into a pollutant-specific AQHI. National mean PM$_{2.5}$ fell
\textbf{55.06 $\to$ 26.50\,$\mu$g\,m$^{-3}$} (2013 $\to$ 2022) then rebounded to 27.46 in
2023; per 10\,$\mu$g\,m$^{-3}$ over three days the associated increases were 0.42\%
all-cause, 0.61\% cardiovascular and 0.84\% respiratory mortality. Note what neither paper
delivers: a \emph{measurement} path to source-resolved toxicity. Both weight after the
fact, from inventories and meta-analysis. For distributed low-cost sensing this is the
open brief --- speciation or source discrimination at network cost, not another mass
channel.
\end{keybox}

\vspace{3mm}

% -------------------------------------------------------------- exec brief
\section*{\sffamily\bfseries\color{Deep}What changed today}
\vspace{-1.5mm}{\color{Amber}\rule{\linewidth}{1.1pt}}\vspace{2.5mm}

\footnotesize
\begin{itemize}
\item \textbf{An unusually modelling-heavy day.} \textit{Exposure assessment \&
modelling} is the largest cluster (5/16) and \textit{Modelling / inventory} the largest
architecture (5/16); ten of sixteen records carry no human health endpoint at all. Only
one record (Wang L et al.) reports an effect estimate with a confidence interval, so
today's forest plot has a single source and is not a synthesis.

\item \textbf{HEI's multinational exposure-assessment initiative reports out.} Boogaard
et al.\ synthesise the whole programme: exposure surfaces that add daily mobility and
outdoor--indoor infiltration correlate $>$0.7 with conventional residential estimates,
and the resulting health effect estimates differ in \emph{magnitude but not direction}.
That is a genuinely useful null for anyone deciding whether a mobility-aware exposure
model is worth its cost --- and it is the strongest argument in today's corpus for
spending effort on validation data rather than on more model structure.

\item \textbf{Sulfate chemistry got two independent corrections, both pointing the same
way.} Liu T et al.\ show aerosol solute strength drives Mn(II)-catalysed sulfate
formation \textbf{more than two orders of magnitude faster} than dilute-cloudwater
parameterisations predict, raising modelled haze-episode sulfate by $>$50\%; Guo et al.\
show that ignoring kinetic isotope effects in source apportionment overstates coal
combustion by 10\% and understates traffic by 8\% in summer. Both say the transition-metal
pathway has been systematically under-credited.

\item \textbf{A cohort result with a real interaction term.} Wang L et al.\ (CHARLS,
$n$=6{,}457, 710 strokes) report \textbf{HR 1.321 (1.222--1.428) per 1\,SD annual
PM$_{2.5}$}, a significant multiplicative interaction with small dense LDL-C
($p$=0.037), and \textbf{HR 1.713 (1.335--2.198)} for the dual-high stratum. The lipid
fraction, not total LDL-C, is what carries the interaction (AUC 0.578 vs.\ 0.534).

\item \textbf{Sub-Saharan Africa is still absent; North Africa is not.} Saad et al.\
report a two-campaign, 24-site port--urban survey at Alexandria --- the first African
record since 29 August. Particulate levels were low across the study area; the port
signal is in CO and NH$_3$, not PM, which the authors state plainly.

\item \textbf{Two records ship on metadata only.} Both indoor-air papers
(\texttt{10.1016/j.buildenv.2026.115193}, \texttt{10.1016/j.indenv.2026.100191})
deposited no abstract to Crossref and are not yet in PubMed, Europe PMC or Semantic
Scholar. They are listed and counted but not summarised, and no numbers are attributed
to them.
\end{itemize}

\clearpage

% ==================================================================== ANALYTICS
\section*{\sffamily\bfseries\color{Deep}Corpus analytics}
\vspace{-1.5mm}{\color{Amber}\rule{\linewidth}{1.1pt}}\vspace{3mm}

\noindent\includegraphics[width=\linewidth]{\FIGDIR/f1_subtopics.png}
\figcap{Subtopic assignment of the 16 in-scope records. Seven of the ten clusters are
occupied; \textit{Neuro / mental health}, \textit{Respiratory \& allergic} and
\textit{Reproductive \& developmental} are all empty today --- the second issue in which
all three are simultaneously vacant, after 5 August. \textit{Exposure assessment \&
modelling} leads with 5, which is the day's defining feature rather than an artefact of
labelling: three of those five (Guo, Chi, Hern\'andez Pardo) are atmospheric-process
papers with no health endpoint at all.}

\vspace{2mm}
\noindent\includegraphics[width=\linewidth]{\FIGDIR/f2_design_endpoint.png}
\figcap{Left --- study architecture. Modelling and inventory work is the largest single
group (5/16), with chamber/laboratory (2), cohort (2) and review (2) behind it; the two
\textit{Metadata only} entries are the abstract-free indoor-air deposits and should be
read as an absence of information, not as a design. Right --- health endpoint by organ
system. \textbf{Ten of sixteen records (62\%) carry no human health endpoint} --- among
the highest shares this month, though short of 20 August's 72\%; the remaining six split
across
musculoskeletal, cardiovascular, respiratory-dosimetry and all-cause mortality.}

\vspace{2mm}
\noindent\includegraphics[width=\linewidth]{\FIGDIR/f3_metric_geography.png}
\figcap{Left --- particle metric. PM$_{2.5}$ alone is the exposure in 7 records and
appears jointly with PM$_{10}$ in 2 more; one record (Garcia Dominguez et al.) is
explicitly ultrafine, and one is bioaerosol/coarse. Right --- geographic provenance.
\textit{Global / multi-region} is the largest bucket at 7, which here reflects the
modelling weight of the day rather than genuine geographic spread --- five of those seven
are simulations or global syntheses with no field site. China contributes 4, Europe 3
(two UK, one Sweden), Southeast Asia 1 (Malaysia) and Middle East \& North Africa 1
(Egypt). \textbf{Sub-Saharan Africa, South Asia and Latin America are all absent.}}

\vspace{2mm}
\noindent\includegraphics[width=\linewidth]{\FIGDIR/f6_lifecourse.png}
\figcap{Life-course windows addressed. Counts are non-exclusive and are deliberately
sparse today: only four records study identifiable human populations. \textbf{No record
in this issue reports an age-stratified estimate}, so the 65+ cell is scored zero even
though both cohorts (UK Biobank, median 13\,y follow-up from baseline age 40--69; CHARLS,
9\,y from age 45+) certainly contain large older strata. Scoring those as 65+ records
would misdescribe what the papers report.}

\vspace{2mm}
\noindent\includegraphics[width=\linewidth]{\FIGDIR/f4_heatmap.png}
\figcap{Subtopic $\times$ study-architecture cross-tabulation. The structure to note is
the near-total separation of the modelling column from any health-endpoint row:
\textit{Exposure assessment \& modelling} and \textit{Sensing, forecasting \&
instrumentation} together hold 8 of 16 records and contribute \emph{zero} effect
estimates. The two cohort records sit alone in \textit{Cardiovascular \& metabolic} and
\textit{Other clinical endpoints}. This is the instrumentation--epidemiology siloing this
watch has flagged repeatedly, and today it is close to complete.}

\vspace{2mm}
\noindent\includegraphics[width=\linewidth]{\FIGDIR/f5_forest.png}
\figcap{\textbf{Both estimates come from a single paper} (Wang L et al., CHARLS) and this
figure is therefore a display of that study's two contrasts, not a synthesis across the
corpus. The per-1\,SD hazard ratio and the dual-high-stratum hazard ratio are not
independent --- the second is a subset contrast of the same cohort --- so the apparent
agreement between them carries no replication weight. No other record in this issue
reported an effect estimate with a confidence interval.}

\clearpage

% ==================================================================== DIGEST
\section*{\sffamily\bfseries\color{Deep}Paper-level digest}
\vspace{-1.5mm}{\color{Amber}\rule{\linewidth}{1.1pt}}\vspace{1mm}

{\fontsize{7.4}{9}\selectfont\color{Slate}Ordered by cluster. Each entry gives the
methodological core and the finding that matters, not an abstract paraphrase.\par}

\band{Exposure assessment \& modelling\hfill 5 records}{Deep}

\paperentry{Li Q et al.\ 2026}
{Health risk of fine PM across China with a pollutant-specific index, 2013--2023}
{Environ Sci Technol}
{A 1\,km daily PM$_{2.5}$ surface for 2013--2023 built by machine learning
(cross-validated $R^2$ \textbf{0.88 spatially, 0.82 temporally}), fused with
meta-analysed short-term effect estimates into an AQHI$_{\mathrm{PM2.5}}$. National mean
fell \textbf{55.06\,$\to$\,26.50\,$\mu$g\,m$^{-3}$} (2013$\to$2022) then rebounded to
27.46 in 2023. Per 10\,$\mu$g\,m$^{-3}$ over 3\,d: \textbf{+0.42\% all-cause, +0.61\%
cardiovascular, +0.84\% respiratory mortality}. Northern and southwestern regions retain
$>$80 medium/high-risk days a year; eliminating them would avert 6{,}985 (6{,}419--7{,}588)
deaths, 7.42\% of the attributable burden. The weak point is the meta-analysed exposure
--- response step, which imports between-study heterogeneity into a spatially resolved
index and then reports it at 1\,km as if the precision were the model's. For sensing
work: the AQHI construction is directly portable to a dense network, but only if the
network can supply the exposure surface it assumes.}
{10.1021/acs.est.6c00144}

\paperentry{Boogaard et al.\ 2026}
{Advancing long-term outdoor air pollution exposure assessment for health studies}
{Curr Environ Health Rep}
{Overview of the completed HEI multinational initiative on novel long-term exposure
assessment across North America and Europe. Exposure predictions from models that account
for daily mobility patterns and outdoor--indoor infiltration correlated
\textbf{$>$0.7} with conventional residential estimates, and the health effect estimates
they produced were similar --- differing in magnitude but \emph{not direction}. The paper
is honest that this is a comparison among models, not a validation against measurement:
representative validation data for long-term exposure simply does not exist, which is the
initiative's own headline gap. Read as the strongest current argument for spending
resources on validation campaigns rather than additional model structure, and as a caution
against over-reading small differences between LUR variants in cohort re-analyses.}
{10.1007/s40572-026-00560-8}

\paperentry{Guo et al.\ 2026b}
{Kinetic isotope fractionation framework for sulfate source apportionment}
{Atmos Chem Phys}
{Seasonal PM$_{2.5}$ observations in Nanjing combined with Bayesian isotope mixing and
process-specific Rayleigh fractionation. Conventional isotope models assume complete
SO$_2$ oxidation and therefore discard the kinetic isotope effects generated during
incomplete processing; correcting for them shows transition-metal-ion-catalysed and
NO$_2$-mediated pathways dominate secondary sulfate, and moves apportionment by
\textbf{$-$10\% coal combustion, $+$8\% traffic} during summer photochemical episodes.
Single-city, single-year, and the Rayleigh treatment assumes a closed reservoir, which is
a strong assumption for an urban boundary layer. The direction of the correction
nevertheless agrees with Liu T et al.\ in the same issue, from wholly independent
evidence.}
{10.5194/acp-26-12231-2026}

\paperentry{Chi et al.\ 2026}
{Heterogeneous nitrosation of amines by N$_2$O$_4$: a missing source of particulate
nitrosamines}
{Atmos Chem Phys}
{Born--Oppenheimer molecular dynamics plus metadynamics on methylamine and
dimethylamine reacting with N$_2$O$_4$ at the air--water interface. Two competing routes:
a \emph{barrierless} direct N-nitrosation yielding nitrosamine cations and NO$_3^-$, and
an amine-mediated N$_2$O$_4$ hydrolysis generating interfacial HONO in
\textbf{2--16\,ps}, which then nitrosates over a 7.65\,kcal\,mol$^{-1}$ barrier at 300\,K.
Purely computational, two amines, idealised interface --- there is no measurement here and
no yield that a chemical transport model can use directly. Its value is as a mechanism
for particulate carcinogen formation that current CTMs omit entirely.}
{10.5194/acp-26-12261-2026}

\paperentry{Hern\'andez Pardo et al.\ 2026}
{Downward transport of tropical upper-tropospheric aerosols}
{Atmos Chem Phys}
{Multi-year global EMAC simulations releasing 19 idealised tracers in the tropical upper
troposphere under resolved advection, parameterised convection, turbulent mixing and
deposition. Age of air at 500\,hPa is $\approx$45\,d for tropics-wide sources but
$>$250\,d for regional continental ones --- transport diagnostics are dominated by source
extent, not by transport physics. Advection dominates; injection height matters more than
particle size in the 20--100\,nm range. Tracers are idealised, carry no mass or
composition, and there is no surface validation, so this constrains \emph{where}
upper-tropospheric particles arrive, not \emph{how much} they contribute to surface PM.
Carried at tier C for that reason.}
{10.5194/acp-26-12275-2026}

\band{Sensing, forecasting \& instrumentation\hfill 3 records}{Amber}

\paperentry{Liu T et al.\ 2026}
{Aerosol solute strength drives TMI-catalysed sulfate aerosol formation}
{Environ Sci Technol}
{Aerosol flow-reactor experiments on atmospherically relevant aqueous aerosols,
decoupling solute strength, interfacial chemistry and Fe(III)--Mn(II) interaction.
Sulfate formation from Mn(II) chemistry runs \textbf{$>$2 orders of magnitude faster}
than models assuming dilute cloudwater; substantial insoluble Fe(III) suppresses the rate
below an ionic-strength threshold of \textbf{16\,mol\,kg$^{-1}$}. Implemented in an
atmospheric model, this raises sulfate abundance during polluted haze episodes by
\textbf{$>$50\%}. The flow reactor cannot reproduce real aerosol organic coatings or
mixing state, so the transfer to ambient particles is an extrapolation --- but this is a
rare case of a laboratory rate measurement closing a known model--measurement gap rather
than opening a new one.}
{10.1021/acs.est.6c09521}

\paperentry{Yohanna \& Lim 2026}
{Regional and seasonal PM$_{2.5}$ and ozone variation across Malaysia, XGBoost + SHAP}
{Chemosphere}
{Daily DOE network observations 2018--2023 for five Malaysian regions, augmented with
ERA5 boundary-layer height, solar radiation and cloud cover; XGBoost against a multiple
linear regression baseline, interpreted with SHAP. \textbf{PM$_{2.5}$ $R^2$=0.576,
RMSE 4.87\,$\mu$g\,m$^{-3}$}; O$_3$ $R^2$=0.543, RMSE 3.4\,ppb --- rising to $R^2$=0.85
only when persistence is included, which is the number to be sceptical of, since a
persistence term will carry most of the skill in any daily-mean series. Kruskal--Wallis
confirms regional ($\chi^2$=3239.5) and seasonal ($\chi^2$=477.8) structure. SHAP
attributes PM variability to combustion tracers and meteorology, ozone chiefly to
humidity and temperature. Useful as a tropical-monsoon reference for anyone tuning a
regional model; not a calibration result.}
{10.1016/j.chemosphere.2026.145074}

\paperentry{Saad et al.\ 2026}
{Short-term air pollutant monitoring at the Alexandria port--urban exposure zone}
{Environ Monit Assess}
{Two monitoring campaigns, May--October 2022, across 11 in-port and 13 surrounding urban
sites, using a portable multi-gas monitor and particulate sensors, with hourly
meteorology from Ras El-Tin (campaign-mean wind 3.0--7.2\,m\,s$^{-1}$, northerly to
north-westerly). Fuel and bulk-fertiliser berths recorded significantly higher CO and
NH$_3$ than urban sites; SO$_2$ and traffic pollutants outside the port dominate the
urban background. \textbf{Particulate levels were low across the whole study area} ---
the authors say so, and the paper is correspondingly framed as screening-level spatial
evidence, not exposure assessment. No sensor model, no reference co-location and no
uncertainty budget are reported, which caps what the PM numbers can support. Carried
mainly because African monitoring records are rare in this corpus.}
{10.1007/s10661-026-15828-7}

\band{Occupational \& indoor exposure\hfill 3 records}{Clay}

\paperentry{Illa et al.\ 2026}
{Exposure pathways when psychoactive drugs are smoked indoors}
{Environ Sci Technol}
{Mechanistic indoor chemical-fate and exposure model across seven compounds spanning
orders of magnitude in volatility and hydrophobicity, with scenarios varying smoker
retention rate and smoking frequency, resolving first-, second- and third-hand pathways
and inhalation, mouthing-ingestion and dermal routes. Dominant pathway is strongly
chemical-specific: volatile compounds (nicotine) go by inhalation; intermediate
hydrophobicity by dermal absorption; involatile hydrophobic compounds
($\Delta^9$-THC) receive meaningful contributions from all three. For non-users the
SHS/THS balance is set jointly by partitioning and by the ratio of active-smoking to
post-smoking time. Entirely modelled --- no measurements, no validation --- so the
useful output is the ranking of pathways, not the doses.}
{10.1021/acs.est.6c02953}

\paperentry{Jin et al.\ 2026b}
{Air change rate and air purifier type for bioaerosol control during dental ultrasonic
scaling}
{Build Environ}
{\textbf{Metadata only.} The Crossref deposit carries no abstract and the record is not
yet in PubMed, Europe PMC or Semantic Scholar. Listed and counted for completeness; no
finding, effect size or method detail is attributed to it, and it is on the abstract
retry list.}
{10.1016/j.buildenv.2026.115193}

\paperentry{Wood et al.\ 2026b}
{Influence of ventilation rate on the potential for indoor particulate matter exposure}
{Indoor Environments}
{\textbf{Metadata only.} Same situation --- no abstract in the Crossref deposit and no
other index carries it yet. Flagged here because the title is squarely in scope for
indoor PM sensing; it is on the abstract retry list and will be summarised if the
abstract appears.}
{10.1016/j.indenv.2026.100191}

\band{Burden, policy \& mitigation\hfill 2 records}{Coral}

\paperentry{Song et al.\ 2026}
{Shifting PM$_{2.5}$ control from mass to toxicity for health equity worldwide}
{Environ Sci Technol}
{Perspective synthesising toxicological, epidemiological and modelling evidence against
the uniform-toxicity assumption underlying mass-based standards. The empirical hook is
that age-standardised attributable death rates diverged worldwide over the past decade
despite falling mass. Incomplete-combustion sources --- residential solid fuel and
wildfire above all --- carry the highest toxicity per unit mass from primary emission
through secondary ageing. Toxicity-adjusted exposure is substantially elevated in South
Asia, Southeast Asia, sub-Saharan Africa, Andean Latin America and China, and the
\textbf{8.2\% of the world that is low-income bears 24.0\% of the cumulative
toxicity-adjusted exposure burden}. It is a perspective, not a quantitative assessment:
the toxicity weights are inherited from a heterogeneous literature and the equity
arithmetic is only as good as they are.}
{10.1021/acs.est.6c05282}

\paperentry{Lim et al.\ 2026}
{Public knowledge, attitudes and practices on air quality in the ten most populous
countries}
{Nat Commun}
{Quota sample of \textbf{10{,}618 respondents across 1{,}372 localities} in Bangladesh,
Brazil, China, Indonesia, India, Mexico, Nigeria, Pakistan, Russia and the United States,
linked to modelled PM$_{2.5}$. Exposure is \emph{inversely} associated with knowledge
about air quality but \emph{positively} associated with protective practices, attitudes
and concern --- people in dirtier air act more and know less. Interactions with age, sex,
education, residence, political ideology, marital, employment and smoking status are all
significant. Quota sampling is not probability sampling, so national point estimates
should not be read as prevalence; the cross-country contrasts are the durable part. A
useful anchor for anyone arguing that a dense sensor network has a communication function
as well as a measurement one.}
{10.1038/s41467-026-75908-7}

\clearpage

\band{Cardiovascular \& metabolic\hfill 1 record}{Teal}

\paperentry{Wang L et al.\ 2026}
{Joint exposure to small dense LDL-C and long-term PM$_{2.5}$ and incident stroke}
{Lipids}
{CHARLS 2011--2020, \textbf{6{,}457} stroke-free adults, \textbf{710} incident strokes
over a median 9\,y; PM$_{2.5}$ from the ChinaHighPM$_{2.5}$ product; Cox models with
$k$-means trajectory clustering and restricted cubic splines. Per 1\,SD:
\textbf{sd-LDL-C HR 1.116 (1.018--1.224)}, \textbf{PM$_{2.5}$ HR 1.321 (1.222--1.428)},
with a significant multiplicative interaction ($p$=0.037) and \textbf{HR 1.713
(1.335--2.198)} for the dual-high stratum against dual-low. A sustained high sd-LDL-C
trajectory added 36\% (1.361, 1.086--1.706). sd-LDL-C outperformed conventional LDL-C for
discrimination (AUC 0.578 vs.\ 0.534, $p<$0.001) --- though both are poor in absolute
terms, which the abstract does not say. sd-LDL-C is \emph{estimated}, not measured, and
PM$_{2.5}$ is a modelled annual mean at residence, so the interaction is between two
proxies and is the result most likely to attenuate under better measurement.}
{10.1002/lipd.70072}

\band{Other clinical endpoints\hfill 1 record}{Slate}

\paperentry{Feng et al.\ 2026}
{Air pollutant mixtures, plasma metabolic and proteomic signatures, and incident arthritis}
{Joint Bone Spine}
{UK Biobank, \textbf{401{,}676} arthritis-free at baseline, \textbf{78{,}188} incident
cases over a median 13\,y. Mixture exposure quantified by weighted quantile sum
regression; metabolic and proteomic signatures derived by elastic net; Cox models plus
generalised propensity scores; causal mediation for the signatures. The dominant mixture
contributors are subtype-specific --- NO$_x$/NO$_2$ for rheumatoid arthritis, PM for
osteoarthritis, NO$_x$ and PM for gout, PM for psoriatic arthritis --- and the omics
signatures mediated \textbf{4.92--35.05\%} of the association. The abstract reports no
hazard ratios or intervals for the WQS index itself, which is why nothing from this paper
appears in the forest plot. WQS indices are also not comparable across studies, and the
mediation proportions rest on a no-unmeasured-confounding assumption that plasma omics
cannot support. The subtype-resolved contributor pattern is the transferable result.}
{10.1016/j.jbspin.2026.106129}

\band{Mechanistic toxicology\hfill 1 record}{Violet}

\paperentry{Garcia Dominguez et al.\ 2026}
{Ultrafine particle deposition and clearance in an avian lung}
{Environ Sci Technol}
{Zebra finches exposed to 50 or 100\,nm model particles at 5 or 25\,$^\circ$C, all within
environmentally relevant ranges, with deposition quantified in lung and --- in a subset ---
heart, liver and red blood cells. Lung deposition was highest for the smallest particles,
consistent with diffusional deposition. Deposited \emph{fraction} was unchanged by
temperature but deposited \emph{dose rate} rose in the cold, driven entirely by elevated
minute ventilation: a clean demonstration that thermoregulatory demand, not respiratory
physiology, sets avian dose. \textbf{44\% of the initially deposited 100\,nm burden
remained at two weeks}, i.e.\ clearance is slow. Model particles, one species, small
subset for extrapulmonary organs --- but this is the only record in today's corpus that
measures ultrafine dose in a living respiratory system, and the avian lung's
cross-current architecture makes it a genuinely different test case from rodent work.}
{10.1021/acs.est.6c05957}

\clearpage

% ==================================================================== TABLE
\section*{\sffamily\bfseries\color{Deep}Corpus register}
\vspace{-1.5mm}{\color{Amber}\rule{\linewidth}{1.1pt}}\vspace{2.5mm}

{\fontsize{6.7}{8.4}\selectfont
\setlength{\tabcolsep}{2.4pt}
\begin{longtable}{@{}>{\RaggedRight}p{25mm}>{\RaggedRight}p{41mm}>{\RaggedRight}p{27mm}%
>{\RaggedRight}p{23mm}>{\RaggedRight}p{18mm}>{\RaggedRight}p{15mm}@{}}
\toprule
\textsf{\textbf{First author}} & \textsf{\textbf{Focus}} & \textsf{\textbf{Journal}} &
\textsf{\textbf{Design}} & \textsf{\textbf{Metric}} & \textsf{\textbf{Region}}\\
\midrule
\endfirsthead
\toprule
\textsf{\textbf{First author}} & \textsf{\textbf{Focus}} & \textsf{\textbf{Journal}} &
\textsf{\textbf{Design}} & \textsf{\textbf{Metric}} & \textsf{\textbf{Region}}\\
\midrule
\endhead
\multicolumn{6}{@{}l@{}}{\cellcolor{Mist}\textsf{\textbf{\textcolor{Deep}{Neuro / mental health}}}}\\[0.6mm]
\href{https://doi.org/10.1001/jamanetworkopen.2026.25015}{Dales} & Wildfire-sourced PM$_{2.5}$ and ED visits for migraine and other primary headache & \textit{JAMA Netw Open} & Case-crossover & PM$_{2.5}$ (wildfire) & Canada\\
\href{https://doi.org/10.1093/ageing/afag215}{Glover} & Umbrella meta-review of urban environmental exposures and cognitive health & \textit{Age Ageing} & Umbrella review / meta-analysis & PM$_{2.5}$ & Global\\
\href{https://doi.org/10.1016/j.nsa.2026.107019}{Goossens} & In utero air pollution exposure and pre-/postnatal brain development & \textit{Neurosci Appl} & Narrative review & PM$_{2.5}$ / PM$_{10}$ / PM$_{\mathrm{coarse}}$ & Global\\
\href{https://doi.org/10.1016/j.isci.2026.116617}{Ma} & PM$_{2.5}$-induced neurotoxicity and depression: exposure risks and mechanisms & \textit{iScience} & Narrative review & PM$_{2.5}$ + components & Global\\
\href{https://doi.org/10.3760/cma.j.cn112150-20250804-00752}{Pan} & Maternal PM exposure and offspring neurodevelopment: effects and mechanisms & \textit{Zhonghua Yu Fang Yi Xue Za Zhi} & Narrative review & PM (bulk) & China\\
\addlinespace[0.9mm]
\multicolumn{6}{@{}l@{}}{\cellcolor{Mist}\textsf{\textbf{\textcolor{Teal}{Cardiovascular \& metabolic}}}}\\[0.6mm]
\href{https://doi.org/10.3760/cma.j.cn112338-20251120-00843}{Huang} & PM$_{2.5}$ and dyslipidaemia in rural Chinese children; dietary effect modification & \textit{Zhonghua Liu Xing Bing Xue Za Zhi} & Cross-sectional (multistage) & PM$_{2.5}$ & China\\
\href{https://doi.org/10.32394/pe/218966}{Kazmierska} & Air pollution and cardiovascular disease: review of current data & \textit{Przegl Epidemiol} & Narrative review & PM$_{2.5}$ / PM$_{10}$ & Global\\
\href{https://doi.org/10.1007/s00125-026-06814-2}{Xi} & Air pollution, built environment and temperature vs. prediabetes phenotypes & \textit{Diabetologia} & Prospective cohort & PM$_{2.5}$ & Germany\\
\href{https://doi.org/10.1161/JAHA.125.046917}{Yang} & Short-term PM$_{2.5}$ and hospitalisation/ED use in atrial fibrillation & \textit{J Am Heart Assoc} & Case-crossover & PM$_{2.5}$ (wildfire-influenced) & USA\\
\addlinespace[0.9mm]
\multicolumn{6}{@{}l@{}}{\cellcolor{Mist}\textsf{\textbf{\textcolor{Sky}{Reproductive \& developmental}}}}\\[0.6mm]
\href{https://doi.org/10.1016/j.jes.2025.10.043}{Chen} & Prenatal air pollution, fetal growth trajectory phenotypes and child neurodevelopment & \textit{J Environ Sci (China)} & Prospective cohort & PM$_{2.5}$ / PM$_{10}$ & China\\
\href{https://doi.org/10.1016/j.ecoenv.2026.120566}{Huang} & Joint prenatal bisphenol + air-pollutant mixture exposure and preschool IQ & \textit{Ecotoxicol Environ Saf} & Prospective cohort (mixtures) & PM$_{2.5}$ / PM$_{10}$ & China\\
\href{https://doi.org/10.1007/s40572-026-00552-8}{Peterson} & Ambient air pollution, PAHs and pubertal development: critical review & \textit{Curr Environ Health Rep} & Critical review & PM$_{2.5}$ components & Global\\
\href{https://doi.org/10.1016/j.envres.2026.125313}{Xie} & Long-term PM$_{2.5}$ and timing of pubertal development in boys and girls & \textit{Environ Res} & Prospective cohort & PM$_{2.5}$ & China\\
\addlinespace[0.9mm]
\multicolumn{6}{@{}l@{}}{\cellcolor{Mist}\textsf{\textbf{\textcolor{Sage}{Respiratory \& allergic}}}}\\[0.6mm]
\href{https://doi.org/10.1007/s00484-026-03289-z}{Bayar Muluk} & Climate change and otolaryngology: airway disease impacts and emission reduction & \textit{Int J Biometeorol} & Narrative review & PM$_{2.5}$ & Global\\
\href{https://doi.org/10.4103/mgr.MEDGASRES-D-25-00164}{Liao} & Short-term air pollution, respiratory symptoms and lung function in asthma & \textit{Med Gas Res} & Panel study & PM$_{2.5}$ / PM$_{10}$ & China\\
\href{https://doi.org/10.3390/epidemiologia7040097}{Sanchez de Prada} & Ambient PM and asthma-coded admissions in children 0-3 y, Spain & \textit{Epidemiologia} & Ecological panel & PM$_{2.5}$ / PM$_{10}$ & Spain\\
\href{https://doi.org/10.1016/j.jhazmat.2026.143076}{Yang} & Allergenic vs. toxic risk of airborne seed hairs compared with ambient PM$_{2.5}$ & \textit{J Hazard Mater} & Field measurement & PM$_{2.5}$ + bioaerosol & China\\
\addlinespace[0.9mm]
\multicolumn{6}{@{}l@{}}{\cellcolor{Mist}\textsf{\textbf{\textcolor{Violet}{Mechanistic toxicology}}}}\\[0.6mm]
\href{https://doi.org/10.1016/j.jes.2025.12.074}{Ding} & PM$_{2.5}$ promotes vascular calcification via miR-27b-5p/HES1 signalling & \textit{J Environ Sci (China)} & In vitro + organoid + murine & PM$_{2.5}$ & China\\
\href{https://doi.org/10.1007/s00360-026-01698-5}{Maraschi} & Osmoregulatory disruption and metal bioaccumulation from settleable atmospheric PM & \textit{J Comp Physiol B} & Controlled exposure (ecotox) & Settleable APM & Brazil\\
\href{https://doi.org/10.1016/j.jes.2025.11.018}{Wu} & Macrophage NLRP3 inflammasome in real-ambient PM-induced hepatic glucose dysregulation & \textit{J Environ Sci (China)} & Real-ambient murine + co-culture & Real-ambient PM & China\\
\addlinespace[0.9mm]
\multicolumn{6}{@{}l@{}}{\cellcolor{Mist}\textsf{\textbf{\textcolor{Clay}{Occupational \& indoor}}}}\\[0.6mm]
\href{https://doi.org/10.3389/fmedt.2026.1796828}{Chen} & Nested suction instrument to cut intraoperative surgical-smoke PM$_{2.5}$ & \textit{Front Med Technol} & Bench / device experiment & PM$_{2.5}$ (surgical smoke) & China\\
\href{https://doi.org/10.1093/annweh/wxag055}{Cole} & Determinants of respirable crystalline silica exposure in Australian tunnelling & \textit{Ann Work Expo Health} & Observational exposure study & Respirable dust / RCS & Australia\\
\href{https://doi.org/10.1016/j.jhazmat.2026.143081}{Ding} & Stage-resolved hair Zn-isotope dynamics as a firefighting combustion-exposure tracer & \textit{J Hazard Mater} & Repeated-measures biomarker & Combustion-derived PM & China\\
\href{https://doi.org/10.3390/jox16040119}{Ramsden} & Xenobiotic hazards in aircraft cabin air & \textit{J Xenobiot} & Review + risk economics & Ultrafine particles & Global\\
\addlinespace[0.9mm]
\multicolumn{6}{@{}l@{}}{\cellcolor{Mist}\textsf{\textbf{\textcolor{Amber}{Sensing, forecasting \& instrumentation}}}}\\[0.6mm]
\href{https://doi.org/10.1016/j.jenvman.2026.130515}{Shi} & Adaptive spatiotemporal graph transformer for multi-site multi-step PM$_{2.5}$ forecasting & \textit{J Environ Manage} & Deep-learning model & PM$_{2.5}$ & China / India\\
\href{https://doi.org/10.3390/bios16070380}{Tsuchiya} & Seasonal and indoor-outdoor characteristics of fluorescent aerosol particles in offices & \textit{Biosensors} & Real-time sensor campaign & Bioaerosol / coarse & Japan\\
\addlinespace[0.9mm]
\multicolumn{6}{@{}l@{}}{\cellcolor{Mist}\textsf{\textbf{\textcolor{Amber}{Exposure assessment \& modelling}}}}\\[0.6mm]
\href{https://doi.org/10.1029/2025GH001621}{Solorzano-Araque} & Multiscale exposure modelling: LOTOS-EUROS + mobility surveys + particle toxicity & \textit{GeoHealth} & CTM + mobility integration & PM$_{2.5}$ / PM$_{10}$ & Colombia\\
\addlinespace[0.9mm]
\multicolumn{6}{@{}l@{}}{\cellcolor{Mist}\textsf{\textbf{\textcolor{Coral}{Burden, policy \& mitigation}}}}\\[0.6mm]
\href{https://doi.org/10.1002/hsr2.72886}{Karamzad} & Ischaemic heart disease burden in MENA 1990-2021; PM attributable fraction 28.2\% & \textit{Health Sci Rep} & GBD secondary analysis & Ambient PM & MENA\\
\href{https://doi.org/10.1371/journal.pone.0347989}{Rees} & HEPA filtration to reduce respiratory infection in care homes & \textit{PLoS One} & RCT process evaluation & Filtered PM & UK\\
\href{https://doi.org/10.1016/j.jes.2026.01.009}{Shan} & Policy-driven emission reductions and post-pandemic rebound, high-resolution inventories & \textit{J Environ Sci (China)} & Emission inventory & PM emissions & China\\
\addlinespace[0.9mm]
\multicolumn{6}{@{}l@{}}{\cellcolor{Mist}\textsf{\textbf{\textcolor{Slate}{Other clinical endpoints}}}}\\[0.6mm]
\href{https://doi.org/10.3390/dj14070440}{Lupita} & PM$_{2.5}$, preventive-care capacity and caries/edentulism burden across the EU & \textit{Dent J} & Ecological panel & PM$_{2.5}$ & EU-27\\
\href{https://doi.org/10.1016/j.jdin.2026.06.008}{Park} & Long-term PM exposure and risk of androgenetic alopecia & \textit{JAAD Int} & Nationwide cohort & PM$_{2.5}$ / PM$_{10}$ & South Korea\\
\href{https://doi.org/10.3171/2026.3.PEDS2615}{Sun} & Environmental exposure and area deprivation vs. paediatric hydrocephalus outcomes & \textit{J Neurosurg Pediatr} & Registry cohort + geospatial & PM$_{2.5}$ (satellite-derived) & USA\\
\addlinespace[0.9mm]
\bottomrule
\end{longtable}}


\vspace{4mm}
\begin{keybox}
{\sffamily\bfseries\footnotesize\color{Deep}Provenance and method}\\[1.2mm]
{\fontsize{7.2}{8.8}\selectfont
\textbf{Window: 31 August 2026 alone}, contiguous with the 30 August issue; no gap. Sources
queried on 1 September 2026, before 22:00 local, so 31 August is the newest buildable date
under the run-timing rule. \textbf{PubMed} \texttt{[EDAT]} 2026/08/31 --- health axis
\textbf{11} via the connector; the instrumentation axis had to be split in two because the
connector caps boolean operators at 20, and both surviving hits (\texttt{42674653},
\texttt{42673782}) were already in the health set, so the sensing axis added \textbf{0}
unique records. \textbf{Europe PMC} \texttt{CREATION\_DATE} 2026-08-31 --- 1 hit
(\texttt{PMID 42669197}, World Trade Center dust), \textbf{already shipped on 30 August};
both the pmid and doi legs of \texttt{seen.json} caught it. \textbf{Crossref by-ISSN}
across the tracked non-PubMed journals --- \textbf{35} deposits with a 31 August
\texttt{created} date, of which 8 are carried and 27 screened out (3 retraction notices,
the rest PFAS, nanoplastics, battery recycling, MOF catalysis, GNSS meteorology and
gas-phase ozone work). \textbf{arXiv} relevance query returned 60 entries, \textbf{none}
published on or after 26 August. \textbf{OpenAlex remains dead} --- HTTP 403 paid-plan
block --- and is routed through \textbf{Consensus}, which returned 3 hits on the
calibration/drift axis: \textbf{all three} resolved to DOIs already in \texttt{seen.json}
(\texttt{10.1088/1361-6501/ae6298}, 29 Jul; \texttt{10.3390/atmos17040335}, 2 Aug;
\texttt{10.5194/amt-19-603-2026}, 1 Aug), so Consensus contributed no backfill this run.
\textbf{47 unique records considered; 1 dropped as already seen; 16 in scope; 31 logged to
\texttt{rejected.jsonl} with a reason.}
\emph{Two records are listed on metadata only} --- \texttt{10.1016/j.buildenv.2026.115193}
and \texttt{10.1016/j.indenv.2026.100191} deposited no abstract to Crossref and are absent
from PubMed, Europe PMC and Semantic Scholar; they are not summarised and no numbers are
attributed to them.
Subtopic, design, particle-metric and geography labels are assigned from title, abstract,
keywords and MeSH terms; assignment is single-label by dominant contribution and is a
judgement, not a controlled vocabulary. First-author surnames are read from PubMed
\texttt{efetch} or the Crossref deposit, never inferred.
\textbf{Trial watch}, window 31 August 2026, ClinicalTrials.gov API v2 on
\texttt{LastUpdatePostDate}: 1 hit, 1 kept --- \texttt{NCT07797413} (Savannah Environmental
Nutrition Study, a food-based antioxidant intervention near the Port of Savannah, $n$=20,
not yet recruiting, feasibility endpoints with FeNO and exhaled-breath-condensate
H$_2$O$_2$ as secondary markers). \texttt{analyze\_endpoints} was run; the record is
surfaced in the monthly rollup, not here.
Content derived from PubMed and Crossref metadata; abstracts remain the copyright of their
respective publishers.\par}
\end{keybox}

\end{document}
